\chapter{The system under test}\label{ch:system}

\section{One turn}

\Cref{fig:turn} follows a customer's message through the copilot. The path in grey is deterministic and runs on
every turn; the teal boxes are the only places a language model may act, and each has a deterministic fallback.

\begin{figure}[htbp]
  \centering
  \resizebox{\linewidth}{!}{% One conversation turn through BETA AID: the deterministic path (ink) and the places a language model may act (teal).
% Kept at the top of tex/ so the number lint, which reads the prose, does not read TikZ coordinates.
\begin{tikzpicture}[
  font=\footnotesize,
  stage/.style={draw=ink, rounded corners=2pt, fill=paper, align=center, minimum height=9mm, text width=23mm, inner sep=2pt},
  model/.style={draw=teal, rounded corners=2pt, fill=teal!8, align=center, text width=23mm, inner sep=2pt, minimum height=9mm},
  stop/.style={draw=brick, rounded corners=2pt, fill=brick!8, align=center, text width=23mm, inner sep=2pt, minimum height=9mm},
  flow/.style={-{Stealth[length=2mm]}, ink, thick},
  side/.style={-{Stealth[length=2mm]}, teal, dashed},
  exit/.style={-{Stealth[length=2mm]}, brick},
  node distance=5mm and 6mm]
  \node[stage] (id) {identity\\ signed session};
  \node[stage, right=of id] (gw) {gateway\\ mask, screen, language};
  \node[stage, right=of gw] (intent) {intent model\\ confident?};
  \node[stage, right=of intent] (card) {card resolution\\ from the session};
  \node[stage, below=12mm of card] (policy) {policy\\ autonomy, vetoes};
  \node[stage, left=of policy] (tools) {verified tools\\ read, confirm, act};
  \node[stage, left=of tools] (reply) {template reply\\ tool facts only};
  \node[stage, left=of reply] (audit) {trace and\\ hash-chained audit};
  \node[model, above=8mm of intent] (classify) {model labels\\ an unsure message};
  \node[model, below=8mm of reply] (rephrase) {model rephrases;\\ grounding check};
  \node[model, below=8mm of tools] (kb) {knowledge base;\\ model answers, cited};
  \node[stop, above=8mm of gw] (refuse) {refused before\\ any model};
  \node[stop, below=8mm of policy] (handoff) {hand-off packet\\ to a person};
  \draw[flow] (id) -- (gw);
  \draw[flow] (gw) -- (intent);
  \draw[flow] (intent) -- (card);
  \draw[flow] (card) -- (policy);
  \draw[flow] (policy) -- (tools);
  \draw[flow] (tools) -- (reply);
  \draw[flow] (reply) -- (audit);
  \draw[side] (intent) -- (classify);
  \draw[side] (reply) -- (rephrase);
  \draw[side] (policy.south west) -- (kb.north east);
  \draw[exit] (gw) -- (refuse);
  \draw[exit] (policy) -- (handoff);
\end{tikzpicture}
}
  \caption{One conversation turn. Grey: the deterministic path. Teal: where a language model may act, each with a
  fallback that needs no model. Brick: where a turn stops before any model or any tool acts.}
  \label{fig:turn}
\end{figure}

\begin{description}[style=nextline]
  \item[Identity.] A signed session token, with an expiry, names the customer. Every tool call takes the customer
  from the verified session, never from the message, so no wording can bring another customer's data into scope.
  Unblocking a card needs a second factor (a step-up) on top of the session.
  \item[Gateway.] Before anything reads the message, card numbers that pass the Luhn check are masked down to the final
  digits a customer uses to recognise a card, prompt-injection and cross-customer probes are flagged and refused, and the language is detected.
  A refused message never reaches a model.
  \item[Intent model.] A character n-gram logistic regression labels the message with one of
  \est{sys.intents}{value} intents (\cref{ch:intent}). When its confidence is below the tuned threshold, a language
  model, if one is configured, labels the message instead; otherwise the copilot asks the customer to rephrase.
  \item[Card resolution.] When the intent needs a card, the copilot resolves which one from the customer's own
  cards, and asks when the message is ambiguous.
  \item[Policy.] A versioned policy file assigns each intent an autonomy level. \est{sys.intents-a0}{value} intents
  are answered from verified facts (\ident{A0}); \est{sys.intents-a2}{value} are card actions that need the
  customer's explicit confirmation (\ident{A2}), \est{sys.intents-step-up}{value} of them also a step-up; and
  \est{sys.intents-a3}{value} go to a person with a hand-off packet (\ident{A3}). \est{sys.vetoes}{value}
  card-level vetoes, checked once the card is known, override the intent: a fraud flag sends the case to the fraud
  queue, a closed card cannot be unblocked, an expired card is replaced rather than unblocked.
  \item[Tools.] Reads come from a read-only snapshot filtered by the session's customer. An action is prepared as a
  confirmation token bound to the action, customer, card and expiry; executed only once, after the customer
  confirms; and reported as done only when a fresh read shows the new state.
  \item[Reply.] A template states the facts the tools returned. A language model may reword it; a grounding check
  rejects any rewording that adds or changes a number, date, amount or card suffix, and the template is used
  instead.
  \item[Knowledge.] A general question goes to the governed knowledge base (\cref{ch:retrieval}). When no
  document scores above the relevance threshold, the copilot asks the customer to rephrase. When the best document
  is internal, only a person may explain it, and the case is handed off. When it is public, a model, if configured,
  writes an answer from the retrieved passages that must cite one of them and contain no number the cited passages
  lack; an answer that fails this check is not shown, and the case is handed off. Without a model, the copilot points
  the customer to the public document.
  \item[Trace and audit.] Every turn records its stages, the policy rule that decided it and the model calls, in a
  hash-chained audit log whose records hold masked text only.
\end{description}

\section{Where a model may act, and where it may not}

The model's jobs are bounded by what it receives and what is done with its output. When it labels an unsure
message it receives the masked message and returns one label from a closed list, which the policy then treats like
any other label: a wrong label for a card action still needs the customer's confirmation, and a wrong hand-off label
only routes the case to a person. When it rewords or answers, its text is checked against the facts or the cited
passages before a customer sees it. It never receives another customer's data, never calls a tool, and never decides
an outcome.

The copilot's model timeout is \est{sys.model-timeout}{value}\,s: a call that has not finished by then is abandoned
and the deterministic path decides. \Cref{ch:screening} measures how often hosted models exceed it.

\section{The governed knowledge base}

The knowledge base holds \est{sys.kb-documents}{value} document versions (\cref{tab:kb-documents}):
\est{sys.kb-approved}{value} approved, \est{sys.kb-draft}{value} draft, \est{sys.kb-superseded}{value} superseded
and \est{sys.kb-retired}{value} retired; \est{sys.kb-public}{value} are classified public. Only approved versions
inside their effective window can be retrieved; drafts, superseded, retired and expired versions stay in the
registry for audit but are never served. The documents are synthetic policies of the demonstration bank and short
summaries of public regulations.

\generated{tables/kb-documents}
